\subsection{Experimental Method Characterization}
\label{sec:method-characterization}

To ensure scientific rigor and reproducibility, this research adopts the Controlled Experiment method. According to \citeonline{wohlin_2012}, experimentation in software engineering provides a systematic, disciplined, quantifiable, and controlled way of evaluating technical solutions. 

The design follows the framework established by \citeonline{basili_1986}, defined as follows:
\begin{itemize}
    \item \textbf{Goal:} Analyze the observability architectures.
    \item \textbf{Purpose:} Evaluate diagnostic efficiency and resource overhead.
    \item \textbf{Perspective:} Researcher and System Operator.
    \item \textbf{Context:} A controlled Kubernetes laboratory environment running a microservices workload under simulated load.
\end{itemize}

The choice of this method is justified by the need to isolate the ``Observability Architecture'' factor from confounding variables, ensuring that observed differences in metrics are indeed caused by the architecture and not by external noise \cite{wohlin_2012}.